% ----------------------------------------------------------------------
%  1.5. Физические свойства кристаллов
% ----------------------------------------------------------------------
\begin{orientir}
Кристаллофизика изучает, как симметрия кристалла ограничивает его
свойства. Два главных инструмента: \vazhno{принципы симметрии Неймана и
Кюри} (1.5.1) и \vazhno{тензорное описание свойств} (1.5.2). Применив их,
узнаем, сколько независимых констант нужно для описания проводимости
(1.5.3) и упругости (1.5.4) кристалла каждой сингонии, а в 1.5.5~—
как это число находится общим методом теории представлений групп.
\end{orientir}

Из-за однородности, анизотропии и симметрии в кристаллах возможны
явления, невозможные в изотропных телах (пьезоэффект, двупреломление,
анизотропия упругости и проводимости). Математический аппарат
кристаллофизики~— тензорное исчисление и теория групп.

\subsection{Принципы симметрии в кристаллофизике}\label{ssec:symmetry-principle}

Для макроскопических свойств кристалл~— однородная среда, поэтому
различия между простыми и винтовыми осями, зеркальными плоскостями и
плоскостями скольжения несущественны. Симметрию макроскопических
свойств определяет не пространственная, а \emph{точечная} группа
кристалла.

\begin{glavnoe}[Принцип Неймана]
Элементы симметрии любого физического свойства кристалла включают
элементы симметрии его точечной группы:
\begin{equation}\label{eq:neumann}
G_\text{свойства}\supseteq G_\text{кристалла} .
\end{equation}
Симметрия свойства \emph{не ниже} симметрии кристалла (может быть и
выше). Точечная группа кристалла~— общая часть (пересечение) групп
симметрии всех его свойств.
\end{glavnoe}

Например, электропроводность кубического кристалла изотропна (симметрия
шара $\infty\infty m$~— намного выше $m\br3m$), а модуль Юнга того же
кристалла анизотропен (рис.~\ref{fig:young-polar}) и имеет ровно
кубическую симметрию.

\begin{glavnoe}[Принцип суперпозиции Кюри]
Кристалл во внешнем поле сохраняет только те элементы симметрии, которые
общи для кристалла без воздействия и для воздействия без кристалла (при
их заданной взаимной ориентации):
\begin{equation}\label{eq:curie}
G_\text{кр+возд}=G_\text{кр}\cap G_\text{возд} .
\end{equation}
\end{glavnoe}

Симметрию воздействий задают предельные группы (п.~\ref{ssec:curie}), а
искать общие элементы удобно по стереографическим проекциям
(рис.~\ref{fig:curie-principle}).
\begin{itemize}
\item \emph{Одноосное растяжение} имеет симметрию $\infty/mmm$
(неподвижный цилиндр). Вдоль оси 4 кристалла $m\br3m$ сохраняются: одна ось
4, четыре поперечные оси 2, четыре продольные и одна поперечная
плоскости, центр~— кристалл становится \emph{тетрагональным} ($4/mmm$).
Вдоль оси 3 остаются ось 3, три поперечные оси 2, три продольные
плоскости и центр~— класс $\br3m$ (\emph{тригональный}).
\item \emph{Электрическое поле}~— полярный вектор, симметрия $\infty m$
(неподвижный конус). Поле вдоль оси 4 кристалла $m\br3m$ даёт класс $4mm$,
вдоль оси 3~— $3m$; центр симметрии исчезает, и кристалл становится
полярным.
\end{itemize}
Обратная задача так не решается: по симметрии результата нельзя судить о
симметрии причин. Обобщённая формулировка Кюри: при наложении явлений
их \emph{диссимметрии} (по Шубникову~— отсутствующие элементы симметрии)
складываются.

\begin{figure}[ht]
\centering
\tikzset{mirror/.style={fmBlue, line width=0.6pt},
  ax4/.style={rectangle, draw=fmRed, fill=fmRed!35, inner sep=0pt, minimum size=5pt},
  ax3/.style={regular polygon, regular polygon sides=3, draw=fmAmber!90!black, fill=fmAmber!40, inner sep=0pt, minimum size=7pt},
  ax2/.style={ellipse, draw=fmGreen!60!black, fill=fmGreen!35, inner sep=0pt, minimum width=7pt, minimum height=4pt}}
\begin{subfigure}[b]{0.32\textwidth}\centering
\begin{tikzpicture}[scale=1.55]
  \draw[mirror, line width=1pt] (0,0) circle (1);
  % создано tools/gen_stereo.py: элементы симметрии m-3m на проекции (001)
\draw[mirror] (0.0000,-1.0000) -- (0.0000,-0.9742) -- (0.0000,-0.9490) -- (0.0000,-0.9244) -- (0.0000,-0.9004) -- (0.0000,-0.8770) -- (0.0000,-0.8541) -- (0.0000,-0.8317) -- (0.0000,-0.8098) -- (0.0000,-0.7883) -- (0.0000,-0.7673) -- (0.0000,-0.7467) -- (0.0000,-0.7265) -- (0.0000,-0.7067) -- (0.0000,-0.6873) -- (0.0000,-0.6682) -- (0.0000,-0.6494) -- (0.0000,-0.6310) -- (0.0000,-0.6128) -- (0.0000,-0.5949) -- (0.0000,-0.5774) -- (0.0000,-0.5600) -- (0.0000,-0.5430) -- (0.0000,-0.5261) -- (0.0000,-0.5095) -- (0.0000,-0.4931) -- (0.0000,-0.4770) -- (0.0000,-0.4610) -- (0.0000,-0.4452) -- (0.0000,-0.4296) -- (0.0000,-0.4142) -- (0.0000,-0.3990) -- (0.0000,-0.3839) -- (0.0000,-0.3689) -- (0.0000,-0.3541) -- (0.0000,-0.3395) -- (0.0000,-0.3249) -- (0.0000,-0.3105) -- (0.0000,-0.2962) -- (0.0000,-0.2820) -- (0.0000,-0.2679) -- (0.0000,-0.2540) -- (0.0000,-0.2401) -- (0.0000,-0.2263) -- (0.0000,-0.2126) -- (0.0000,-0.1989) -- (0.0000,-0.1853) -- (0.0000,-0.1718) -- (0.0000,-0.1584) -- (0.0000,-0.1450) -- (0.0000,-0.1317) -- (0.0000,-0.1184) -- (0.0000,-0.1051) -- (0.0000,-0.0919) -- (0.0000,-0.0787) -- (0.0000,-0.0655) -- (0.0000,-0.0524) -- (0.0000,-0.0393) -- (0.0000,-0.0262) -- (0.0000,-0.0131) -- (0.0000,-0.0000) -- (-0.0000,0.0131) -- (-0.0000,0.0262) -- (-0.0000,0.0393) -- (-0.0000,0.0524) -- (-0.0000,0.0655) -- (-0.0000,0.0787) -- (-0.0000,0.0919) -- (-0.0000,0.1051) -- (-0.0000,0.1184) -- (-0.0000,0.1317) -- (-0.0000,0.1450) -- (-0.0000,0.1584) -- (-0.0000,0.1718) -- (-0.0000,0.1853) -- (-0.0000,0.1989) -- (-0.0000,0.2126) -- (-0.0000,0.2263) -- (-0.0000,0.2401) -- (-0.0000,0.2540) -- (-0.0000,0.2679) -- (-0.0000,0.2820) -- (-0.0000,0.2962) -- (-0.0000,0.3105) -- (-0.0000,0.3249) -- (-0.0000,0.3395) -- (-0.0000,0.3541) -- (-0.0000,0.3689) -- (-0.0000,0.3839) -- (-0.0000,0.3990) -- (-0.0000,0.4142) -- (-0.0000,0.4296) -- (-0.0000,0.4452) -- (-0.0000,0.4610) -- (-0.0000,0.4770) -- (-0.0000,0.4931) -- (-0.0000,0.5095) -- (-0.0000,0.5261) -- (-0.0000,0.5430) -- (-0.0000,0.5600) -- (-0.0000,0.5774) -- (-0.0000,0.5949) -- (-0.0000,0.6128) -- (-0.0000,0.6310) -- (-0.0000,0.6494) -- (-0.0000,0.6682) -- (-0.0000,0.6873) -- (-0.0000,0.7067) -- (-0.0000,0.7265) -- (-0.0000,0.7467) -- (-0.0000,0.7673) -- (-0.0000,0.7883) -- (-0.0000,0.8098) -- (-0.0000,0.8317) -- (-0.0000,0.8541) -- (-0.0000,0.8770) -- (-0.0000,0.9004) -- (-0.0000,0.9244) -- (-0.0000,0.9490) -- (-0.0000,0.9742) -- (-0.0000,1.0000);
\draw[mirror] (1.0000,0.0000) -- (0.9742,0.0000) -- (0.9490,0.0000) -- (0.9244,0.0000) -- (0.9004,0.0000) -- (0.8770,0.0000) -- (0.8541,0.0000) -- (0.8317,0.0000) -- (0.8098,0.0000) -- (0.7883,0.0000) -- (0.7673,0.0000) -- (0.7467,0.0000) -- (0.7265,0.0000) -- (0.7067,0.0000) -- (0.6873,0.0000) -- (0.6682,0.0000) -- (0.6494,0.0000) -- (0.6310,0.0000) -- (0.6128,0.0000) -- (0.5949,0.0000) -- (0.5774,0.0000) -- (0.5600,0.0000) -- (0.5430,0.0000) -- (0.5261,0.0000) -- (0.5095,0.0000) -- (0.4931,0.0000) -- (0.4770,0.0000) -- (0.4610,0.0000) -- (0.4452,0.0000) -- (0.4296,0.0000) -- (0.4142,0.0000) -- (0.3990,0.0000) -- (0.3839,0.0000) -- (0.3689,0.0000) -- (0.3541,0.0000) -- (0.3395,0.0000) -- (0.3249,0.0000) -- (0.3105,0.0000) -- (0.2962,0.0000) -- (0.2820,0.0000) -- (0.2679,0.0000) -- (0.2540,0.0000) -- (0.2401,0.0000) -- (0.2263,0.0000) -- (0.2126,0.0000) -- (0.1989,0.0000) -- (0.1853,0.0000) -- (0.1718,0.0000) -- (0.1584,0.0000) -- (0.1450,0.0000) -- (0.1317,0.0000) -- (0.1184,0.0000) -- (0.1051,0.0000) -- (0.0919,0.0000) -- (0.0787,0.0000) -- (0.0655,0.0000) -- (0.0524,0.0000) -- (0.0393,0.0000) -- (0.0262,0.0000) -- (0.0131,0.0000) -- (0.0000,0.0000) -- (-0.0131,-0.0000) -- (-0.0262,-0.0000) -- (-0.0393,-0.0000) -- (-0.0524,-0.0000) -- (-0.0655,-0.0000) -- (-0.0787,-0.0000) -- (-0.0919,-0.0000) -- (-0.1051,-0.0000) -- (-0.1184,-0.0000) -- (-0.1317,-0.0000) -- (-0.1450,-0.0000) -- (-0.1584,-0.0000) -- (-0.1718,-0.0000) -- (-0.1853,-0.0000) -- (-0.1989,-0.0000) -- (-0.2126,-0.0000) -- (-0.2263,-0.0000) -- (-0.2401,-0.0000) -- (-0.2540,-0.0000) -- (-0.2679,-0.0000) -- (-0.2820,-0.0000) -- (-0.2962,-0.0000) -- (-0.3105,-0.0000) -- (-0.3249,-0.0000) -- (-0.3395,-0.0000) -- (-0.3541,-0.0000) -- (-0.3689,-0.0000) -- (-0.3839,-0.0000) -- (-0.3990,-0.0000) -- (-0.4142,-0.0000) -- (-0.4296,-0.0000) -- (-0.4452,-0.0000) -- (-0.4610,-0.0000) -- (-0.4770,-0.0000) -- (-0.4931,-0.0000) -- (-0.5095,-0.0000) -- (-0.5261,-0.0000) -- (-0.5430,-0.0000) -- (-0.5600,-0.0000) -- (-0.5774,-0.0000) -- (-0.5949,-0.0000) -- (-0.6128,-0.0000) -- (-0.6310,-0.0000) -- (-0.6494,-0.0000) -- (-0.6682,-0.0000) -- (-0.6873,-0.0000) -- (-0.7067,-0.0000) -- (-0.7265,-0.0000) -- (-0.7467,-0.0000) -- (-0.7673,-0.0000) -- (-0.7883,-0.0000) -- (-0.8098,-0.0000) -- (-0.8317,-0.0000) -- (-0.8541,-0.0000) -- (-0.8770,-0.0000) -- (-0.9004,-0.0000) -- (-0.9244,-0.0000) -- (-0.9490,-0.0000) -- (-0.9742,-0.0000) -- (-1.0000,-0.0000);
\draw[mirror] (0.7071,-0.7071) -- (0.6888,-0.6888) -- (0.6710,-0.6710) -- (0.6536,-0.6536) -- (0.6367,-0.6367) -- (0.6201,-0.6201) -- (0.6039,-0.6039) -- (0.5881,-0.5881) -- (0.5726,-0.5726) -- (0.5574,-0.5574) -- (0.5426,-0.5426) -- (0.5280,-0.5280) -- (0.5137,-0.5137) -- (0.4997,-0.4997) -- (0.4860,-0.4860) -- (0.4725,-0.4725) -- (0.4592,-0.4592) -- (0.4462,-0.4462) -- (0.4333,-0.4333) -- (0.4207,-0.4207) -- (0.4082,-0.4082) -- (0.3960,-0.3960) -- (0.3839,-0.3839) -- (0.3720,-0.3720) -- (0.3603,-0.3603) -- (0.3487,-0.3487) -- (0.3373,-0.3373) -- (0.3260,-0.3260) -- (0.3148,-0.3148) -- (0.3038,-0.3038) -- (0.2929,-0.2929) -- (0.2821,-0.2821) -- (0.2714,-0.2714) -- (0.2609,-0.2609) -- (0.2504,-0.2504) -- (0.2400,-0.2400) -- (0.2298,-0.2298) -- (0.2196,-0.2196) -- (0.2095,-0.2095) -- (0.1994,-0.1994) -- (0.1895,-0.1895) -- (0.1796,-0.1796) -- (0.1698,-0.1698) -- (0.1600,-0.1600) -- (0.1503,-0.1503) -- (0.1407,-0.1407) -- (0.1311,-0.1311) -- (0.1215,-0.1215) -- (0.1120,-0.1120) -- (0.1025,-0.1025) -- (0.0931,-0.0931) -- (0.0837,-0.0837) -- (0.0743,-0.0743) -- (0.0650,-0.0650) -- (0.0557,-0.0557) -- (0.0463,-0.0463) -- (0.0371,-0.0371) -- (0.0278,-0.0278) -- (0.0185,-0.0185) -- (0.0093,-0.0093) -- (0.0000,-0.0000) -- (-0.0093,0.0093) -- (-0.0185,0.0185) -- (-0.0278,0.0278) -- (-0.0371,0.0371) -- (-0.0463,0.0463) -- (-0.0557,0.0557) -- (-0.0650,0.0650) -- (-0.0743,0.0743) -- (-0.0837,0.0837) -- (-0.0931,0.0931) -- (-0.1025,0.1025) -- (-0.1120,0.1120) -- (-0.1215,0.1215) -- (-0.1311,0.1311) -- (-0.1407,0.1407) -- (-0.1503,0.1503) -- (-0.1600,0.1600) -- (-0.1698,0.1698) -- (-0.1796,0.1796) -- (-0.1895,0.1895) -- (-0.1994,0.1994) -- (-0.2095,0.2095) -- (-0.2196,0.2196) -- (-0.2298,0.2298) -- (-0.2400,0.2400) -- (-0.2504,0.2504) -- (-0.2609,0.2609) -- (-0.2714,0.2714) -- (-0.2821,0.2821) -- (-0.2929,0.2929) -- (-0.3038,0.3038) -- (-0.3148,0.3148) -- (-0.3260,0.3260) -- (-0.3373,0.3373) -- (-0.3487,0.3487) -- (-0.3603,0.3603) -- (-0.3720,0.3720) -- (-0.3839,0.3839) -- (-0.3960,0.3960) -- (-0.4082,0.4082) -- (-0.4207,0.4207) -- (-0.4333,0.4333) -- (-0.4462,0.4462) -- (-0.4592,0.4592) -- (-0.4725,0.4725) -- (-0.4860,0.4860) -- (-0.4997,0.4997) -- (-0.5137,0.5137) -- (-0.5280,0.5280) -- (-0.5426,0.5426) -- (-0.5574,0.5574) -- (-0.5726,0.5726) -- (-0.5881,0.5881) -- (-0.6039,0.6039) -- (-0.6201,0.6201) -- (-0.6367,0.6367) -- (-0.6536,0.6536) -- (-0.6710,0.6710) -- (-0.6888,0.6888) -- (-0.7071,0.7071);
\draw[mirror] (-0.7071,-0.7071) -- (-0.6888,-0.6888) -- (-0.6710,-0.6710) -- (-0.6536,-0.6536) -- (-0.6367,-0.6367) -- (-0.6201,-0.6201) -- (-0.6039,-0.6039) -- (-0.5881,-0.5881) -- (-0.5726,-0.5726) -- (-0.5574,-0.5574) -- (-0.5426,-0.5426) -- (-0.5280,-0.5280) -- (-0.5137,-0.5137) -- (-0.4997,-0.4997) -- (-0.4860,-0.4860) -- (-0.4725,-0.4725) -- (-0.4592,-0.4592) -- (-0.4462,-0.4462) -- (-0.4333,-0.4333) -- (-0.4207,-0.4207) -- (-0.4082,-0.4082) -- (-0.3960,-0.3960) -- (-0.3839,-0.3839) -- (-0.3720,-0.3720) -- (-0.3603,-0.3603) -- (-0.3487,-0.3487) -- (-0.3373,-0.3373) -- (-0.3260,-0.3260) -- (-0.3148,-0.3148) -- (-0.3038,-0.3038) -- (-0.2929,-0.2929) -- (-0.2821,-0.2821) -- (-0.2714,-0.2714) -- (-0.2609,-0.2609) -- (-0.2504,-0.2504) -- (-0.2400,-0.2400) -- (-0.2298,-0.2298) -- (-0.2196,-0.2196) -- (-0.2095,-0.2095) -- (-0.1994,-0.1994) -- (-0.1895,-0.1895) -- (-0.1796,-0.1796) -- (-0.1698,-0.1698) -- (-0.1600,-0.1600) -- (-0.1503,-0.1503) -- (-0.1407,-0.1407) -- (-0.1311,-0.1311) -- (-0.1215,-0.1215) -- (-0.1120,-0.1120) -- (-0.1025,-0.1025) -- (-0.0931,-0.0931) -- (-0.0837,-0.0837) -- (-0.0743,-0.0743) -- (-0.0650,-0.0650) -- (-0.0557,-0.0557) -- (-0.0463,-0.0463) -- (-0.0371,-0.0371) -- (-0.0278,-0.0278) -- (-0.0185,-0.0185) -- (-0.0093,-0.0093) -- (-0.0000,-0.0000) -- (0.0093,0.0093) -- (0.0185,0.0185) -- (0.0278,0.0278) -- (0.0371,0.0371) -- (0.0463,0.0463) -- (0.0557,0.0557) -- (0.0650,0.0650) -- (0.0743,0.0743) -- (0.0837,0.0837) -- (0.0931,0.0931) -- (0.1025,0.1025) -- (0.1120,0.1120) -- (0.1215,0.1215) -- (0.1311,0.1311) -- (0.1407,0.1407) -- (0.1503,0.1503) -- (0.1600,0.1600) -- (0.1698,0.1698) -- (0.1796,0.1796) -- (0.1895,0.1895) -- (0.1994,0.1994) -- (0.2095,0.2095) -- (0.2196,0.2196) -- (0.2298,0.2298) -- (0.2400,0.2400) -- (0.2504,0.2504) -- (0.2609,0.2609) -- (0.2714,0.2714) -- (0.2821,0.2821) -- (0.2929,0.2929) -- (0.3038,0.3038) -- (0.3148,0.3148) -- (0.3260,0.3260) -- (0.3373,0.3373) -- (0.3487,0.3487) -- (0.3603,0.3603) -- (0.3720,0.3720) -- (0.3839,0.3839) -- (0.3960,0.3960) -- (0.4082,0.4082) -- (0.4207,0.4207) -- (0.4333,0.4333) -- (0.4462,0.4462) -- (0.4592,0.4592) -- (0.4725,0.4725) -- (0.4860,0.4860) -- (0.4997,0.4997) -- (0.5137,0.5137) -- (0.5280,0.5280) -- (0.5426,0.5426) -- (0.5574,0.5574) -- (0.5726,0.5726) -- (0.5881,0.5881) -- (0.6039,0.6039) -- (0.6201,0.6201) -- (0.6367,0.6367) -- (0.6536,0.6536) -- (0.6710,0.6710) -- (0.6888,0.6888) -- (0.7071,0.7071);
\draw[mirror] (0.0000,-1.0000) -- (-0.0182,-0.9815) -- (-0.0357,-0.9630) -- (-0.0526,-0.9445) -- (-0.0688,-0.9261) -- (-0.0845,-0.9077) -- (-0.0996,-0.8893) -- (-0.1142,-0.8710) -- (-0.1282,-0.8528) -- (-0.1417,-0.8346) -- (-0.1547,-0.8165) -- (-0.1672,-0.7985) -- (-0.1793,-0.7805) -- (-0.1910,-0.7626) -- (-0.2022,-0.7448) -- (-0.2130,-0.7271) -- (-0.2234,-0.7095) -- (-0.2334,-0.6919) -- (-0.2430,-0.6745) -- (-0.2523,-0.6571) -- (-0.2612,-0.6398) -- (-0.2698,-0.6226) -- (-0.2780,-0.6055) -- (-0.2860,-0.5884) -- (-0.2936,-0.5715) -- (-0.3009,-0.5546) -- (-0.3080,-0.5378) -- (-0.3147,-0.5211) -- (-0.3212,-0.5045) -- (-0.3274,-0.4879) -- (-0.3333,-0.4714) -- (-0.3390,-0.4550) -- (-0.3445,-0.4386) -- (-0.3497,-0.4224) -- (-0.3546,-0.4061) -- (-0.3594,-0.3900) -- (-0.3639,-0.3739) -- (-0.3682,-0.3579) -- (-0.3723,-0.3419) -- (-0.3761,-0.3260) -- (-0.3798,-0.3101) -- (-0.3833,-0.2943) -- (-0.3865,-0.2785) -- (-0.3896,-0.2628) -- (-0.3925,-0.2471) -- (-0.3951,-0.2315) -- (-0.3976,-0.2159) -- (-0.4000,-0.2003) -- (-0.4021,-0.1848) -- (-0.4040,-0.1693) -- (-0.4058,-0.1538) -- (-0.4074,-0.1383) -- (-0.4089,-0.1229) -- (-0.4101,-0.1075) -- (-0.4112,-0.0921) -- (-0.4121,-0.0767) -- (-0.4129,-0.0614) -- (-0.4135,-0.0460) -- (-0.4139,-0.0307) -- (-0.4141,-0.0153) -- (-0.4142,-0.0000) -- (-0.4141,0.0153) -- (-0.4139,0.0307) -- (-0.4135,0.0460) -- (-0.4129,0.0614) -- (-0.4121,0.0767) -- (-0.4112,0.0921) -- (-0.4101,0.1075) -- (-0.4089,0.1229) -- (-0.4074,0.1383) -- (-0.4058,0.1538) -- (-0.4040,0.1693) -- (-0.4021,0.1848) -- (-0.4000,0.2003) -- (-0.3976,0.2159) -- (-0.3951,0.2315) -- (-0.3925,0.2471) -- (-0.3896,0.2628) -- (-0.3865,0.2785) -- (-0.3833,0.2943) -- (-0.3798,0.3101) -- (-0.3761,0.3260) -- (-0.3723,0.3419) -- (-0.3682,0.3579) -- (-0.3639,0.3739) -- (-0.3594,0.3900) -- (-0.3546,0.4061) -- (-0.3497,0.4224) -- (-0.3445,0.4386) -- (-0.3390,0.4550) -- (-0.3333,0.4714) -- (-0.3274,0.4879) -- (-0.3212,0.5045) -- (-0.3147,0.5211) -- (-0.3080,0.5378) -- (-0.3009,0.5546) -- (-0.2936,0.5715) -- (-0.2860,0.5884) -- (-0.2780,0.6055) -- (-0.2698,0.6226) -- (-0.2612,0.6398) -- (-0.2523,0.6571) -- (-0.2430,0.6745) -- (-0.2334,0.6919) -- (-0.2234,0.7095) -- (-0.2130,0.7271) -- (-0.2022,0.7448) -- (-0.1910,0.7626) -- (-0.1793,0.7805) -- (-0.1672,0.7985) -- (-0.1547,0.8165) -- (-0.1417,0.8346) -- (-0.1282,0.8528) -- (-0.1142,0.8710) -- (-0.0996,0.8893) -- (-0.0845,0.9077) -- (-0.0688,0.9261) -- (-0.0526,0.9445) -- (-0.0357,0.9630) -- (-0.0182,0.9815) -- (-0.0000,1.0000);
\draw[mirror] (0.0000,-1.0000) -- (0.0182,-0.9815) -- (0.0357,-0.9630) -- (0.0526,-0.9445) -- (0.0688,-0.9261) -- (0.0845,-0.9077) -- (0.0996,-0.8893) -- (0.1142,-0.8710) -- (0.1282,-0.8528) -- (0.1417,-0.8346) -- (0.1547,-0.8165) -- (0.1672,-0.7985) -- (0.1793,-0.7805) -- (0.1910,-0.7626) -- (0.2022,-0.7448) -- (0.2130,-0.7271) -- (0.2234,-0.7095) -- (0.2334,-0.6919) -- (0.2430,-0.6745) -- (0.2523,-0.6571) -- (0.2612,-0.6398) -- (0.2698,-0.6226) -- (0.2780,-0.6055) -- (0.2860,-0.5884) -- (0.2936,-0.5715) -- (0.3009,-0.5546) -- (0.3080,-0.5378) -- (0.3147,-0.5211) -- (0.3212,-0.5045) -- (0.3274,-0.4879) -- (0.3333,-0.4714) -- (0.3390,-0.4550) -- (0.3445,-0.4386) -- (0.3497,-0.4224) -- (0.3546,-0.4061) -- (0.3594,-0.3900) -- (0.3639,-0.3739) -- (0.3682,-0.3579) -- (0.3723,-0.3419) -- (0.3761,-0.3260) -- (0.3798,-0.3101) -- (0.3833,-0.2943) -- (0.3865,-0.2785) -- (0.3896,-0.2628) -- (0.3925,-0.2471) -- (0.3951,-0.2315) -- (0.3976,-0.2159) -- (0.4000,-0.2003) -- (0.4021,-0.1848) -- (0.4040,-0.1693) -- (0.4058,-0.1538) -- (0.4074,-0.1383) -- (0.4089,-0.1229) -- (0.4101,-0.1075) -- (0.4112,-0.0921) -- (0.4121,-0.0767) -- (0.4129,-0.0614) -- (0.4135,-0.0460) -- (0.4139,-0.0307) -- (0.4141,-0.0153) -- (0.4142,-0.0000) -- (0.4141,0.0153) -- (0.4139,0.0307) -- (0.4135,0.0460) -- (0.4129,0.0614) -- (0.4121,0.0767) -- (0.4112,0.0921) -- (0.4101,0.1075) -- (0.4089,0.1229) -- (0.4074,0.1383) -- (0.4058,0.1538) -- (0.4040,0.1693) -- (0.4021,0.1848) -- (0.4000,0.2003) -- (0.3976,0.2159) -- (0.3951,0.2315) -- (0.3925,0.2471) -- (0.3896,0.2628) -- (0.3865,0.2785) -- (0.3833,0.2943) -- (0.3798,0.3101) -- (0.3761,0.3260) -- (0.3723,0.3419) -- (0.3682,0.3579) -- (0.3639,0.3739) -- (0.3594,0.3900) -- (0.3546,0.4061) -- (0.3497,0.4224) -- (0.3445,0.4386) -- (0.3390,0.4550) -- (0.3333,0.4714) -- (0.3274,0.4879) -- (0.3212,0.5045) -- (0.3147,0.5211) -- (0.3080,0.5378) -- (0.3009,0.5546) -- (0.2936,0.5715) -- (0.2860,0.5884) -- (0.2780,0.6055) -- (0.2698,0.6226) -- (0.2612,0.6398) -- (0.2523,0.6571) -- (0.2430,0.6745) -- (0.2334,0.6919) -- (0.2234,0.7095) -- (0.2130,0.7271) -- (0.2022,0.7448) -- (0.1910,0.7626) -- (0.1793,0.7805) -- (0.1672,0.7985) -- (0.1547,0.8165) -- (0.1417,0.8346) -- (0.1282,0.8528) -- (0.1142,0.8710) -- (0.0996,0.8893) -- (0.0845,0.9077) -- (0.0688,0.9261) -- (0.0526,0.9445) -- (0.0357,0.9630) -- (0.0182,0.9815) -- (0.0000,1.0000);
\draw[mirror] (1.0000,0.0000) -- (0.9815,-0.0182) -- (0.9630,-0.0357) -- (0.9445,-0.0526) -- (0.9261,-0.0688) -- (0.9077,-0.0845) -- (0.8893,-0.0996) -- (0.8710,-0.1142) -- (0.8528,-0.1282) -- (0.8346,-0.1417) -- (0.8165,-0.1547) -- (0.7985,-0.1672) -- (0.7805,-0.1793) -- (0.7626,-0.1910) -- (0.7448,-0.2022) -- (0.7271,-0.2130) -- (0.7095,-0.2234) -- (0.6919,-0.2334) -- (0.6745,-0.2430) -- (0.6571,-0.2523) -- (0.6398,-0.2612) -- (0.6226,-0.2698) -- (0.6055,-0.2780) -- (0.5884,-0.2860) -- (0.5715,-0.2936) -- (0.5546,-0.3009) -- (0.5378,-0.3080) -- (0.5211,-0.3147) -- (0.5045,-0.3212) -- (0.4879,-0.3274) -- (0.4714,-0.3333) -- (0.4550,-0.3390) -- (0.4386,-0.3445) -- (0.4224,-0.3497) -- (0.4061,-0.3546) -- (0.3900,-0.3594) -- (0.3739,-0.3639) -- (0.3579,-0.3682) -- (0.3419,-0.3723) -- (0.3260,-0.3761) -- (0.3101,-0.3798) -- (0.2943,-0.3833) -- (0.2785,-0.3865) -- (0.2628,-0.3896) -- (0.2471,-0.3925) -- (0.2315,-0.3951) -- (0.2159,-0.3976) -- (0.2003,-0.4000) -- (0.1848,-0.4021) -- (0.1693,-0.4040) -- (0.1538,-0.4058) -- (0.1383,-0.4074) -- (0.1229,-0.4089) -- (0.1075,-0.4101) -- (0.0921,-0.4112) -- (0.0767,-0.4121) -- (0.0614,-0.4129) -- (0.0460,-0.4135) -- (0.0307,-0.4139) -- (0.0153,-0.4141) -- (0.0000,-0.4142) -- (-0.0153,-0.4141) -- (-0.0307,-0.4139) -- (-0.0460,-0.4135) -- (-0.0614,-0.4129) -- (-0.0767,-0.4121) -- (-0.0921,-0.4112) -- (-0.1075,-0.4101) -- (-0.1229,-0.4089) -- (-0.1383,-0.4074) -- (-0.1538,-0.4058) -- (-0.1693,-0.4040) -- (-0.1848,-0.4021) -- (-0.2003,-0.4000) -- (-0.2159,-0.3976) -- (-0.2315,-0.3951) -- (-0.2471,-0.3925) -- (-0.2628,-0.3896) -- (-0.2785,-0.3865) -- (-0.2943,-0.3833) -- (-0.3101,-0.3798) -- (-0.3260,-0.3761) -- (-0.3419,-0.3723) -- (-0.3579,-0.3682) -- (-0.3739,-0.3639) -- (-0.3900,-0.3594) -- (-0.4061,-0.3546) -- (-0.4224,-0.3497) -- (-0.4386,-0.3445) -- (-0.4550,-0.3390) -- (-0.4714,-0.3333) -- (-0.4879,-0.3274) -- (-0.5045,-0.3212) -- (-0.5211,-0.3147) -- (-0.5378,-0.3080) -- (-0.5546,-0.3009) -- (-0.5715,-0.2936) -- (-0.5884,-0.2860) -- (-0.6055,-0.2780) -- (-0.6226,-0.2698) -- (-0.6398,-0.2612) -- (-0.6571,-0.2523) -- (-0.6745,-0.2430) -- (-0.6919,-0.2334) -- (-0.7095,-0.2234) -- (-0.7271,-0.2130) -- (-0.7448,-0.2022) -- (-0.7626,-0.1910) -- (-0.7805,-0.1793) -- (-0.7985,-0.1672) -- (-0.8165,-0.1547) -- (-0.8346,-0.1417) -- (-0.8528,-0.1282) -- (-0.8710,-0.1142) -- (-0.8893,-0.0996) -- (-0.9077,-0.0845) -- (-0.9261,-0.0688) -- (-0.9445,-0.0526) -- (-0.9630,-0.0357) -- (-0.9815,-0.0182) -- (-1.0000,-0.0000);
\draw[mirror] (1.0000,0.0000) -- (0.9815,0.0182) -- (0.9630,0.0357) -- (0.9445,0.0526) -- (0.9261,0.0688) -- (0.9077,0.0845) -- (0.8893,0.0996) -- (0.8710,0.1142) -- (0.8528,0.1282) -- (0.8346,0.1417) -- (0.8165,0.1547) -- (0.7985,0.1672) -- (0.7805,0.1793) -- (0.7626,0.1910) -- (0.7448,0.2022) -- (0.7271,0.2130) -- (0.7095,0.2234) -- (0.6919,0.2334) -- (0.6745,0.2430) -- (0.6571,0.2523) -- (0.6398,0.2612) -- (0.6226,0.2698) -- (0.6055,0.2780) -- (0.5884,0.2860) -- (0.5715,0.2936) -- (0.5546,0.3009) -- (0.5378,0.3080) -- (0.5211,0.3147) -- (0.5045,0.3212) -- (0.4879,0.3274) -- (0.4714,0.3333) -- (0.4550,0.3390) -- (0.4386,0.3445) -- (0.4224,0.3497) -- (0.4061,0.3546) -- (0.3900,0.3594) -- (0.3739,0.3639) -- (0.3579,0.3682) -- (0.3419,0.3723) -- (0.3260,0.3761) -- (0.3101,0.3798) -- (0.2943,0.3833) -- (0.2785,0.3865) -- (0.2628,0.3896) -- (0.2471,0.3925) -- (0.2315,0.3951) -- (0.2159,0.3976) -- (0.2003,0.4000) -- (0.1848,0.4021) -- (0.1693,0.4040) -- (0.1538,0.4058) -- (0.1383,0.4074) -- (0.1229,0.4089) -- (0.1075,0.4101) -- (0.0921,0.4112) -- (0.0767,0.4121) -- (0.0614,0.4129) -- (0.0460,0.4135) -- (0.0307,0.4139) -- (0.0153,0.4141) -- (0.0000,0.4142) -- (-0.0153,0.4141) -- (-0.0307,0.4139) -- (-0.0460,0.4135) -- (-0.0614,0.4129) -- (-0.0767,0.4121) -- (-0.0921,0.4112) -- (-0.1075,0.4101) -- (-0.1229,0.4089) -- (-0.1383,0.4074) -- (-0.1538,0.4058) -- (-0.1693,0.4040) -- (-0.1848,0.4021) -- (-0.2003,0.4000) -- (-0.2159,0.3976) -- (-0.2315,0.3951) -- (-0.2471,0.3925) -- (-0.2628,0.3896) -- (-0.2785,0.3865) -- (-0.2943,0.3833) -- (-0.3101,0.3798) -- (-0.3260,0.3761) -- (-0.3419,0.3723) -- (-0.3579,0.3682) -- (-0.3739,0.3639) -- (-0.3900,0.3594) -- (-0.4061,0.3546) -- (-0.4224,0.3497) -- (-0.4386,0.3445) -- (-0.4550,0.3390) -- (-0.4714,0.3333) -- (-0.4879,0.3274) -- (-0.5045,0.3212) -- (-0.5211,0.3147) -- (-0.5378,0.3080) -- (-0.5546,0.3009) -- (-0.5715,0.2936) -- (-0.5884,0.2860) -- (-0.6055,0.2780) -- (-0.6226,0.2698) -- (-0.6398,0.2612) -- (-0.6571,0.2523) -- (-0.6745,0.2430) -- (-0.6919,0.2334) -- (-0.7095,0.2234) -- (-0.7271,0.2130) -- (-0.7448,0.2022) -- (-0.7626,0.1910) -- (-0.7805,0.1793) -- (-0.7985,0.1672) -- (-0.8165,0.1547) -- (-0.8346,0.1417) -- (-0.8528,0.1282) -- (-0.8710,0.1142) -- (-0.8893,0.0996) -- (-0.9077,0.0845) -- (-0.9261,0.0688) -- (-0.9445,0.0526) -- (-0.9630,0.0357) -- (-0.9815,0.0182) -- (-1.0000,0.0000);
\node[ax4] at (1.0000,0.0000) {};
\node[ax4] at (-1.0000,0.0000) {};
\node[ax4] at (0.0000,0.0000) {};
\node[ax4] at (0.0000,1.0000) {};
\node[ax4] at (0.0000,-1.0000) {};
\node[ax2, rotate=0] at (0.4142,0.0000) {};
\node[ax2, rotate=180] at (-0.4142,0.0000) {};
\node[ax2, rotate=45] at (0.7071,0.7071) {};
\node[ax2, rotate=-45] at (0.7071,-0.7071) {};
\node[ax2, rotate=135] at (-0.7071,0.7071) {};
\node[ax2, rotate=-135] at (-0.7071,-0.7071) {};
\node[ax2, rotate=90] at (0.0000,0.4142) {};
\node[ax2, rotate=-90] at (0.0000,-0.4142) {};
\node[ax3] at (0.3660,0.3660) {};
\node[ax3] at (0.3660,-0.3660) {};
\node[ax3] at (-0.3660,0.3660) {};
\node[ax3] at (-0.3660,-0.3660) {};

\end{tikzpicture}
\caption{$m\br3m$}
\end{subfigure}\hfill
\begin{subfigure}[b]{0.32\textwidth}\centering
\begin{tikzpicture}[scale=1.55]
  \draw[mirror, line width=1pt] (0,0) circle (1);
  \foreach \a in {0,45,90,135}{\draw[mirror] (\a:1) -- (\a+180:1);
    \node[ax2, rotate=\a] at (\a:1) {}; \node[ax2, rotate=\a] at (\a+180:1) {};}
  \node[ax4] at (0,0) {};
\end{tikzpicture}
\caption{$4/mmm$ (растяжение)}
\end{subfigure}\hfill
\begin{subfigure}[b]{0.32\textwidth}\centering
\begin{tikzpicture}[scale=1.55]
  \draw[atomLine, line width=0.6pt] (0,0) circle (1);
  \foreach \a in {0,45,90,135}{\draw[mirror] (\a:1) -- (\a+180:1);}
  \node[ax4] at (0,0) {};
\end{tikzpicture}
\caption{$4mm$ (поле $\vb E$)}
\end{subfigure}
\caption{Принцип суперпозиции Кюри: кубический кристалл $m\br3m$
(\textit{а}) при одноосном растяжении вдоль $[001]$ (перпендикулярно
плоскости рисунка) сохраняет элементы класса $4/mmm$ (\textit{б}), а в
электрическом поле вдоль $[001]$~— класса $4mm$ (\textit{в}): исчезают
горизонтальная плоскость (основной круг теперь тонкий) и оси 2}
\label{fig:curie-principle}
\end{figure}

\subsection{Тензорное описание физических свойств}\label{ssec:tensors}

\paragraph{Кристаллофизическая система координат}
Свойства описывают в правой прямоугольной системе $x_1x_2x_3$. В
кубической, тетрагональной и ромбической сингониях её оси совпадают с
кристаллографическими; в остальных её ориентируют по стандартным
правилам (табл.~\ref{tab:physaxes}).

\begin{table}[ht]
\centering
\caption{Правила кристаллофизической установки}\label{tab:physaxes}
\small
\begin{tabular}{@{}llll@{}}
\toprule
Сингония & Ось $x_1$ & Ось $x_2$ & Ось $x_3$\\
\midrule
Триклинная & \multicolumn{2}{l}{в плоскости, перпендикулярной $[001]$} & $[001]$\\
Моноклинная & в плоскости $(100)$ & $[010]$ & $[001]$\\
Ромбическая, тетрагон., кубическая & $[100]$ & $[010]$ & $[001]$\\
Тригон. и гексагональная & $[2\br1\br10]$ & $[01\br10]$ & $[0001]$\\
\bottomrule
\end{tabular}
\end{table}

\paragraph{Свойство как связь воздействия и отклика}
Физическое свойство проявляется как коэффициент в соотношении
\begin{equation}\label{eq:property}
\text{явление (отклик)}=\text{свойство}\times\text{воздействие} .
\end{equation}
Ранг тензора свойства определяется рангами воздействия и отклика:
\begin{itemize}
\item скаляр $\to$ скаляр: свойство~— \emph{скаляр} (тензор нулевого
ранга: плотность, теплоёмкость, энтропия);
\item скаляр $\to$ вектор: \emph{вектор} (тензор 1-го ранга). Пример~—
\emph{пироэффект}: при однородном нагреве (скаляр) кристалл поляризуется
($\vb P$~— вектор); возможен только в 10 полярных классах;
\item вектор $\to$ вектор: \emph{тензор 2-го ранга}: электро- и
теплопроводность, удельное электро- и тепловое сопротивление,
диэлектрическая и магнитная проницаемость и восприимчивость~— все они
\emph{симметричные} тензоры 2-го ранга;
\item тензор 2-го ранга $\to$ вектор: \emph{3-й ранг} (пьезоэффект; эффект
Холла: $\rho_{ij}(\vb B)=\rho_{ij}+R_{ijk}B_k$);
\item тензор 2-го ранга $\to$ тензор 2-го ранга: \emph{4-й ранг} (упругость,
магнитострикция).
\end{itemize}

При повороте прямоугольной системы координат (матрица $c$, п.~\ref{ssec:orth})
компоненты тензора 2-го ранга преобразуются как произведения координат:
\begin{equation}\label{eq:tensor2}
A'_{kl}=c_{ki}c_{lj}A_{ij}
\end{equation}
(правило Эйнштейна: по дважды повторяющемуся индексу подразумевается
суммирование от 1 до 3). Тензор \emph{симметричен}, если $A_{ij}=A_{ji}$, и
\emph{антисимметричен}, если $A_{ij}=-A_{ji}$.

\paragraph{Материальные и полевые тензоры}
\emph{Материальные} тензоры описывают свойства кристалла
(проводимость, упругость); их симметрия подчинена принципу Неймана.
\emph{Полевые} тензоры описывают воздействия и отклики (поле $\vb E$,
напряжение, деформация); их можно задать любыми, независимо от симметрии
кристалла. Но отклик на произвольное воздействие уже зависит от симметрии
(например, деформация при заданном напряжении определяется упругими
константами).

\subsection{Симметрия кристалла и тензоры второго ранга}\label{ssec:tensor2-symmetry}

Возьмём тензор электропроводности $\sigma$: $\vb j=\sigma\vb E$, т.\,е.
$j_i=\sigma_{ij}E_j$. При одновременной смене знаков $\vb E$ и $\vb j$
компоненты $\sigma_{ij}$ не меняются, т.\,е. \vazhno{любое свойство,
описываемое тензором 2-го ранга, центросимметрично}~— его симметрия
выше симметрии кристалла по крайней мере на центр инверсии.

Как симметрия «обнуляет» компоненты? Если операция $R$ (матрица)~—
элемент симметрии кристалла, то по принципу Неймана тензор не меняется:
\begin{equation}\label{eq:tensor-invariance}
\sigma'=R\,\sigma\,R^{-1}=\sigma .
\end{equation}
\begin{itemize}
\item \emph{Триклинный} кристалл: ограничений нет, у симметричного тензора
6 независимых компонент.
\item \emph{Моноклинный}, ось 2 вдоль $x_2$: $C_{2y}=\operatorname{diag}(\br1,1,\br1)$,
и
\[
C_{2y}\sigma C_{2y}^{-1}=\begin{pmatrix}\sigma_{11}&-\sigma_{12}&\sigma_{13}\\-\sigma_{12}&\sigma_{22}&-\sigma_{23}\\\sigma_{13}&-\sigma_{23}&\sigma_{33}\end{pmatrix}=\sigma
\ \Rightarrow\ \sigma_{12}=\sigma_{23}=0 ;
\]
остаются 4 компоненты: $\sigma_{11}$, $\sigma_{22}$, $\sigma_{33}$, $\sigma_{13}$.
\item \emph{Ромбический}: оси 2 вдоль $x_1$ и $x_3$ обнуляют и $\sigma_{13}$~— 3
компоненты.
\item \emph{Кубический}: ось 3 вдоль $[111]$ циклически переставляет оси
$x_1\to x_2\to x_3$ и переводит $\operatorname{diag}(\sigma_{11},\sigma_{22},\sigma_{33})$ в
$\operatorname{diag}(\sigma_{33},\sigma_{11},\sigma_{22})$; значит,
$\sigma_{11}=\sigma_{22}=\sigma_{33}=\sigma$~— одна компонента: \vazhno{кубический
кристалл изотропен по всем свойствам 2-го ранга} (но не по упругости!).
\end{itemize}

\begin{table}[ht]
\centering
\caption{Число независимых компонент симметричного тензора 2-го ранга}\label{tab:tensor2}
\small
\begin{tabular}{@{}lcl@{}}
\toprule
Сингония & Число & Независимые компоненты\\
\midrule
Триклинная & 6 & $\sigma_{11}$, $\sigma_{22}$, $\sigma_{33}$, $\sigma_{12}$, $\sigma_{13}$, $\sigma_{23}$\\
Моноклинная & 4 & $\sigma_{11}$, $\sigma_{22}$, $\sigma_{33}$, $\sigma_{13}$\\
Ромбическая & 3 & $\sigma_{11}$, $\sigma_{22}$, $\sigma_{33}$\\
Тетрагональная, тригональная, & 2 & $\sigma_{11}=\sigma_{22}$, $\sigma_{33}$\\
\quad гексагональная & & \\
Кубическая & 1 & $\sigma_{11}=\sigma_{22}=\sigma_{33}=\sigma$\\
\bottomrule
\end{tabular}
\end{table}

\begin{primer}[Проводимость в одноосном кристалле]
У гексагонального (тетрагонального, тригонального) кристалла
$\sigma_{11}=\sigma_{22}=\sigma_\perp$, $\sigma_{33}=\sigma_\parallel$. Для направления под
углом $\theta$ к оси $c$ проводимость (проекция $\vb j$ на $\vb E$)
$\sigma(\theta)=\sigma_\perp\sin^2\theta+\sigma_\parallel\cos^2\theta$. Для цинка, например,
удельное сопротивление вдоль и поперёк оси $c$ различается лишь на
несколько процентов, а у графита (п.~\ref{ssec:elem-types})~— в $10^4$ раз и более.
\end{primer}

\subsection{Упругие свойства кристаллов}\label{ssec:elastic-crystals}

\paragraph{Напряжения и деформации}
Однородное напряжённое состояние задаётся симметричным тензором
напряжений $\sigma_{ij}$: $\sigma_{11}$, $\sigma_{22}$, $\sigma_{33}$~— нормальные
напряжения, остальные~— касательные (сдвиговые). Приведя тензор к
главным осям, получим главные напряжения $\sigma_1$, $\sigma_2$, $\sigma_3$.

Деформацию описывают через смещения $\vb u$ точек тела. Если точки $P$ и
$Q$ отстоят на $\Delta\vb x$, то после деформации разность их смещений
\begin{equation}\label{eq:distortion}
\Delta u_i=u_{ij}\Delta x_j,\qquad u_{ij}=\frac{\partial u_i}{\partial x_j}
\end{equation}
($u_{ij}$~— \term{тензор дисторсии}). Удлинение $e_{11}=\partial u_1/\partial x_1$~—
относительное изменение длины, сдвиг $e_{12}=\partial u_1/\partial x_2=\operatorname{tg}\alpha$~—
смещение одной части тела относительно другой. Но не всякая дисторсия~—
деформация: тело можно просто повернуть. Разложим дисторсию на
симметричную и антисимметричную части (рис.~\ref{fig:distortion}):
\begin{keyeq}\label{eq:strain}
u_{ij}=\underbrace{\tfrac12(u_{ij}+u_{ji})}_{\varepsilon_{ij}}+\underbrace{\tfrac12(u_{ij}-u_{ji})}_{\omega_{ij}} .
\end{keyeq}
Симметричный $\varepsilon_{ij}$~— \term{тензор деформации} (его можно привести к
главным осям и главным деформациям $\varepsilon_1$, $\varepsilon_2$, $\varepsilon_3$);
антисимметричный $\omega_{ij}$~— \term{тензор поворота}, равносильный
аксиальному вектору поворота $\vb\omega=(\omega_1,\omega_2,\omega_3)$, где
$\omega_{12}=-\omega_3$, $\omega_{13}=\omega_2$, $\omega_{23}=-\omega_1$.

\begin{figure}[ht]
\centering
\begin{tikzpicture}[scale=1.25]
  \tikzset{orig/.style={fmGray, dashed, line width=0.5pt}, def/.style={fmBlue, thick, fill=fmBlue!12}}
  \begin{scope}
    \draw[def] (0,0) -- (1.5,0) -- (1.95,1.5) -- (0.45,1.5) -- cycle;
    \draw[orig] (0,0) rectangle (1.5,1.5);
    \draw[fmRed, -{Stealth[length=1.6mm]}] (0,1.5) -- (0.42,1.5);
    \node[font=\small] at (0.75,-0.45) {дисторсия $u_{12}$};
  \end{scope}
  \node[font=\Large] at (2.5,0.75) {$=$};
  \begin{scope}[xshift=3.1cm]
    \draw[def] (0,0) -- (1.5,0.225) -- (1.725,1.725) -- (0.225,1.5) -- cycle;
    \draw[orig] (0,0) rectangle (1.5,1.5);
    \node[font=\small] at (0.75,-0.45) {деформация $\varepsilon_{12}$};
  \end{scope}
  \node[font=\Large] at (5.5,0.75) {$+$};
  \begin{scope}[xshift=6.1cm]
    \draw[def, rotate around={-8.5:(0,0)}] (0,0) rectangle (1.5,1.5);
    \draw[orig] (0,0) rectangle (1.5,1.5);
    \node[font=\small] at (0.75,-0.45) {поворот $\omega$};
  \end{scope}
\end{tikzpicture}
\caption{Простой сдвиг (дисторсия $u_{12}$) есть сумма чистой деформации
(квадрат вытягивается вдоль диагонали, $\varepsilon_{12}=\varepsilon_{21}=u_{12}/2$) и
поворота тела как целого на угол $u_{12}/2$}
\label{fig:distortion}
\end{figure}

\paragraph{Обобщённый закон Гука}
Для изотропного тела $\sigma=c\varepsilon$ ($c$~— жёсткость) или $\varepsilon=s\sigma$
($s=1/c$~— податливость). В кристалле
\begin{keyeq}\label{eq:hooke-gen}
\sigma_{ij}=c_{ijkl}\varepsilon_{kl},\qquad \varepsilon_{ij}=s_{ijkl}\sigma_{kl},
\end{keyeq}
где $c_{ijkl}$ (упругие жёсткости) и $s_{ijkl}$ (податливости)~— тензоры
4-го ранга. Из $3^4=81$ компоненты независимы лишь 21: симметрия тензоров
напряжений и деформаций сводит их число к 36, а существование упругой
энергии ($c_{ijkl}=c_{klij}$)~— к 21. Дальше работает симметрия кристалла.

Удобна \emph{матричная} (фойгтовская) запись: пары индексов заменяют одним:
$11\to1$, $22\to2$, $33\to3$, $23,32\to4$, $13,31\to5$, $12,21\to6$. Тогда
$\sigma_i=c_{ij}\varepsilon_j$, $\varepsilon_i=s_{ij}\sigma_j$ ($i,j=1\ldots6$), а матрицы
$c$ и $s$ взаимно обратны\isprmark. Для кубических и гексагональных кристаллов
\[
c_\text{куб}=\begin{pmatrix}c_{11}&c_{12}&c_{12}&0&0&0\\c_{12}&c_{11}&c_{12}&0&0&0\\c_{12}&c_{12}&c_{11}&0&0&0\\0&0&0&c_{44}&0&0\\0&0&0&0&c_{44}&0\\0&0&0&0&0&c_{44}\end{pmatrix},
\]
\[
c_\text{гекс}=\begin{pmatrix}c_{11}&c_{12}&c_{13}&0&0&0\\c_{12}&c_{11}&c_{13}&0&0&0\\c_{13}&c_{13}&c_{33}&0&0&0\\0&0&0&c_{44}&0&0\\0&0&0&0&c_{44}&0\\0&0&0&0&0&\frac{c_{11}-c_{12}}2\end{pmatrix}
\]
(у матриц $s$~— та же структура, только в последнем элементе
гексагональной матрицы стоит $2(s_{11}-s_{12})$). Итак, кубический кристалл
описывается \vazhno{тремя} упругими константами, гексагональный~—
\vazhno{пятью}, изотропное тело~— \vazhno{двумя} ($s_{11}=1/E$,
$s_{12}=-\nu/E$, $s_{44}=1/G$, причём $G=E/2(1+\nu)$, п.~3.8.1).
\isprtext{В матричной записи сдвиговые компоненты деформации берут
удвоенными: $\varepsilon_4=2\varepsilon_{23}$, $\varepsilon_5=2\varepsilon_{13}$, $\varepsilon_6=2\varepsilon_{12}$
(так, чтобы $c_{44}$ был обычным модулем сдвига); в оригинале этот
множитель не оговорён. Соответственно $s_{44}=4s_{2323}$ и т.\,п.}

\paragraph{Анизотропия упругости кубических кристаллов}
Её меру даёт \term{фактор анизотропии}
\begin{keyeq}\label{eq:zener}
A=\frac{2c_{44}}{c_{11}-c_{12}}=\frac{2(s_{11}-s_{12})}{s_{44}} ;
\end{keyeq}
для изотропного тела $A=1$. Модуль Юнга в направлении с направляющими
косинусами $l_1$, $l_2$, $l_3$:
\begin{keyeq}\label{eq:E-cubic}
\frac1{E}=s_{11}-2\Bigl(s_{11}-s_{12}-\frac{s_{44}}2\Bigr)\bigl(l_1^2l_2^2+l_2^2l_3^2+l_3^2l_1^2\bigr),
\end{keyeq}
для гексагональных кристаллов (угол $\theta$ с осью $c$, $l_3=\cos\theta$)
\begin{equation}\label{eq:E-hex}
\frac1E=(1-l_3^2)^2s_{11}+l_3^4s_{33}+l_3^2(1-l_3^2)(2s_{13}+s_{44}) .
\end{equation}
Сумма $l_1^2l_2^2+l_2^2l_3^2+l_3^2l_1^2$ равна 0 для $\nfam{100}$ и $\frac13$ для
$\nfam{111}$. Поэтому при $A>1$ (тогда $s_{11}-s_{12}-s_{44}/2>0$) самое
жёсткое направление~— $\nfam{111}$, самое податливое~— $\nfam{100}$; при
$A<1$~— наоборот (табл.~\ref{tab:cubic-elastic}, рис.~\ref{fig:young-polar}).
Это определяется величиной $A$, а не типом решётки: среди ОЦК-металлов
есть и $\alpha$-Fe, Ta ($A>1$), и Cr, Mo, V, Nb ($A<1$)\ispr{В оригинале
сказано, что у ГЦК-кристаллов модуль Юнга растёт при переходе от $\nfam{100}$ к
$\nfam{111}$, а у ОЦК~— убывает. Таблица того же оригинала опровергает это:
у ОЦК-железа и тантала $E_{111}>E_{001}$. Направление роста определяется
знаком $A-1$.}. Почти изотропны вольфрам ($A\approx0{,}98$) и алюминий
($A\approx1{,}2$)\ispr{В оригинале наиболее изотропным среди ОЦК-металлов
назван молибден ($A=0{,}91$), хотя у вольфрама $A=0{,}98$.}.

\begin{table}[ht]
\centering
\caption{Упругие постоянные кубических монокристаллов (ГПа)}\label{tab:cubic-elastic}
\small
\begin{tabular}{@{}lcccccccc@{}}
\toprule
Металл & $c_{11}$ & $c_{12}$ & $c_{44}$ & $E_{001}$ & $E_{011}$ & $E_{111}$ & $E_\text{ср}$ & $A$\\
\midrule
Al & 108 & 61 & 28 & 63 & 72 & 76 & 70 & 1,23\\
Cu & 168 & 121 & 75 & 67 & 130 & 191 & 127 & 3,21\\
Ag & 124 & 93 & 46 & 43 & 83 & 120 & 81 & 3,04\\
Au & 186 & 157 & 42 & 43 & 81 & 116 & 78 & 2,87\\
Ni & 246 & 147 & 125 & 130 & 223 & 294 & 214 & 2,54\\
Pb & 50 & 42 & 15 & 11 & 24 & 40 & 24 & 4,12\\
$\alpha$-Fe & 231 & 135 & 116 & 132 & 220 & 283 & 211 & 2,41\\
Ta & 267 & 161 & 82 & 142 & 191 & 215 & 183 & 1,58\\
V & 228 & 119 & 43 & 146 & 124 & 117 & 128 & 0,78\\
Nb & 246 & 139 & 29 & 152 & 92 & 82 & 105 & 0,51\\
Cr & 350 & 67 & 101 & 328 & 266 & 250 & 280 & 0,71\\
Mo & 460 & 176 & 110 & 344 & 326 & 316 & 327 & 0,91\\
W & 501 & 198 & 151 & 420 & 414 & 412 & 415 & 0,98\\
UO$_2$ & 395 & 121 & 64 & 338 & 199 & 175 & 230 & 0,47\\
\bottomrule
\end{tabular}
\end{table}

\begin{table}[ht]
\centering
\caption{Упругие постоянные гексагональных монокристаллов (ГПа)}\label{tab:hex-elastic}
\small
\begin{tabular}{@{}lcccccccc@{}}
\toprule
Металл & $c_{11}$ & $c_{12}$ & $c_{13}$ & $c_{33}$ & $c_{44}$ & $c_{66}$ & $E_\text{ср}$ & $c/a$\\
\midrule
Be & 292 & 27 & 14 & 336 & 162 & 133 & 311 & 1,57\\
Mg & 53 & 26 & 21 & 61 & 16 & 17 & 44 & 1,62\\
Ti & 162 & 92 & 69 & 181 & 47 & 35 & 114 & 1,59\\
$\alpha$-Zr & 144 & 73 & 66 & 165 & 32 & 36 & 97 & 1,59\\
Hf & 181 & 77 & 66 & 197 & 56 & 52 & 143 & 1,58\\
Co & 307 & 165 & 105 & 357 & 76 & 71 & 216 & 1,62\\
Re & 613 & 270 & 206 & 683 & 162 & 171 & 462 & 1,62\\
Zn & 165 & 31 & 50 & 62 & 40 & 67 & 102 & 1,86\\
Cd & 114 & 40 & 60 & 51 & 20 & 38 & 63 & 1,89\\
BeO & 470 & 168 & 119 & 494 & 153 & 151 & 394 & 1,62\\
\bottomrule
\multicolumn{9}{l}{\footnotesize $c_{66}=(c_{11}-c_{12})/2$.}
\end{tabular}
\end{table}

\noindent В гексагональной таблице исправлены опечатки
оригинала\ispr{В табл.~1.8 оригинала для кобальта $c_{33}=36$~ГПа (правильно
около 357~ГПа; при 36~ГПа кристалл был бы вдесятеро мягче вдоль оси $c$,
чем поперёк), для кадмия $c/a=1{,}62$ (в тексте того же оригинала верно:
1,89), для BeO $c/a=1{,}58$ (у BeO со структурой вюрцита $c/a\approx1{,}62$).}.
Обратите внимание на цинк и кадмий с аномальным $c/a$: вдоль оси $c$ они
гораздо мягче ($c_{33}\ll c_{11}$)~— связь между базисными слоями слабее, чем
внутри слоя.

\subsection{Матричное представление группы}\label{ssec:representations}

Самый общий способ выяснить, сколько независимых компонент оставляет
симметрия у тензора любого ранга,~— теория представлений групп.

\begin{opredelenie}[Представление группы]
Сопоставление каждому элементу $R$ группы $G$ матрицы $D(R)$ так, что
произведению элементов отвечает произведение матриц: если $AB=C$, то
$D(A)D(B)=D(C)$. Число строк матриц~— \term{размерность} представления.
\end{opredelenie}

\begin{itemize}
\item Представления $D$ и $D'$ \emph{эквивалентны}, если
$D'(R)=SD(R)S^{-1}$ для всех $R$ с одной и той же матрицей $S$.
\item Представление \emph{приводимо}, если подходящей заменой базиса все
его матрицы одновременно приводятся к блочно-диагональному виду; тогда
оно~— прямая сумма представлений меньшей размерности:
$D(R)=D^{(1)}(R)\oplus D^{(2)}(R)$. Основную роль играют неприводимые
неэквивалентные представления.
\item \emph{Характер} представления~— набор следов $\chi(R)=\operatorname{tr}D(R)$.
Элементы одного класса имеют одинаковый след, поэтому характер~—
функция класса.
\end{itemize}

\begin{primer}[Тензор 2-го ранга при повороте вокруг оси 4]
Пусть ось 4 направлена вдоль $x_1$. Из $\sigma'=C_{4x}\sigma C_{4x}^{-1}$ получаем
$\sigma'_{11}=\sigma_{11}$, $\sigma'_{22}=\sigma_{33}$, $\sigma'_{33}=\sigma_{22}$, $\sigma'_{12}=-\sigma_{13}$,
$\sigma'_{13}=\sigma_{12}$, $\sigma'_{23}=-\sigma_{23}$. Записав шесть компонент столбцом
$(\sigma_{11},\sigma_{22},\sigma_{33},\sigma_{12},\sigma_{13},\sigma_{23})$, получаем
шестимерную матрицу $D(C_{4x})$, которая распадается в прямую сумму
одномерных $(1)$ и $(-1)$ и двумерных
$\bigl(\begin{smallmatrix}0&1\\1&0\end{smallmatrix}\bigr)$,
$\bigl(\begin{smallmatrix}0&-1\\1&0\end{smallmatrix}\bigr)$ блоков. Её след
$\chi=1+0+0+0+0-1=0$, что совпадает с формулой~\eqref{eq:chi-sym} ниже при
$\varphi=90^\circ$.
\end{primer}

\paragraph{Теоремы ортогональности}
Для неприводимых неэквивалентных представлений $D^{(i)}$ размерности $l_i$
группы порядка $h$
\begin{equation}\label{eq:orthogonality}
\sum_{R\in G}D^{(i)}(R)_{\mu\nu}\,D^{(j)}(R)^*_{\alpha\beta}=\frac h{l_i}\,\delta_{ij}\delta_{\mu\alpha}\delta_{\nu\beta}
\end{equation}
(суммирование~— по всем \emph{элементам группы})\ispr{В оригинале
сказано, что суммирование в теореме ортогональности ведётся «по всем
неэквивалентным, неприводимым представлениям»; на самом деле сумма берётся
по элементам группы $R$, а индексы $i$, $j$ нумеруют представления.}. Отсюда
следуют: $\sum_il_i^2=h$; ортогональность характеров
$\sum_kN_k\chi^{(i)}(C_k)\chi^{(j)}(C_k)^*=h\delta_{ij}$ (сумма по классам $C_k$, $N_k$~—
число элементов в классе) и то, что число неприводимых представлений
равно числу классов.

\paragraph{Сколько независимых компонент у тензора}
Компонента, не меняющаяся ни при одной операции группы, отвечает
\emph{единичному} (полносимметричному) представлению с $\chi\equiv1$. Число
таких независимых комбинаций~— кратность единичного представления в
представлении тензора:
\begin{keyeq}\label{eq:n-invariants}
n=\frac1h\sum_kN_k\,\chi(C_k) ,
\end{keyeq}
где $\chi(C_k)$~— след матрицы, по которой преобразуются компоненты тензора,
для операций класса $C_k$. Для поворота на угол $\varphi$ ($\det R=+1$) и
несобственного поворота ($\det R=-1$; поворот на $\varphi$ с отражением в
перпендикулярной плоскости)
\begin{equation}\label{eq:chi-sym}
\chi_\text{сим}=4\cos^2\varphi\pm2\cos\varphi\quad\text{(симметричный тензор 2-го ранга)},
\end{equation}
\begin{equation}\label{eq:chi-gen}
\chi=(\pm1+2\cos\varphi)^2\quad\text{(произвольный тензор 2-го ранга)}
\end{equation}
(знак «$+$» для собственных, «$-$» для несобственных операций).

\begin{primer}[Ромбический кристалл $mmm$]
Элементы группы ($h=8$, каждый~— свой класс): $E$ ($\varphi=0$, $\chi=6$); три
оси 2 ($\varphi=180^\circ$, $\chi=4-2=2$); центр инверсии~— несобственный
поворот на $180^\circ$ ($\chi=4+2=6$); три плоскости~— несобственные повороты на
$0^\circ$ ($\chi=4-2=2$). По~\eqref{eq:n-invariants}
$n=\frac18(6+3\cdot2+6+3\cdot2)=3$: у симметричного тензора 2-го ранга
ромбического кристалла (например, электропроводности $\alpha$-урана) три
независимые компоненты~— как и в табл.~\ref{tab:tensor2}.
\end{primer}

\begin{itogi}
\begin{itemize}
\item Принцип Неймана: $G_\text{свойства}\supseteq G_\text{кристалла}$; принцип Кюри:
$G_\text{кр+возд}=G_\text{кр}\cap G_\text{возд}$.
\item Свойства~— тензоры: скаляры, векторы (пироэффект), 2-го ранга
(проводимость, проницаемость), 3-го (пьезоэффект), 4-го (упругость).
\item Свойства 2-го ранга центросимметричны; число независимых компонент:
6, 4, 3, 2, 1 (от триклинной до кубической сингонии); кубический
кристалл изотропен по ним.
\item Упругость: $\sigma_{ij}=c_{ijkl}\varepsilon_{kl}$; 21 константа в общем
случае, 5~— у гексагональных, 3~— у кубических, 2~— у изотропных;
анизотропия кубических кристаллов $A=2c_{44}/(c_{11}-c_{12})$.
\item Число независимых компонент: $n=\frac1h\sum_kN_k\chi(C_k)$.
\end{itemize}
\end{itogi}

\begin{proverka}
\item Какую симметрию приобретёт кубический кристалл при сжатии вдоль
$[110]$?
\item Почему кубический кристалл не может быть пироэлектриком?
\item Почему электропроводность кубического кристалла изотропна, а
модуль Юнга~— нет?
\item Чему равен модуль Юнга меди вдоль $\nfam{110}$, если
$c_{11}=168$, $c_{12}=121$, $c_{44}=75$~ГПа?
\item Найдите по~\eqref{eq:n-invariants} число независимых компонент
симметричного тензора 2-го ранга для класса $4/m$.
\item Почему у цинка $c_{33}$ много меньше $c_{11}$?
\end{proverka}
